\section{Introduction}
An input/output memory management unit (IOMMU) lets direct memory access (DMA) devices use translated addresses while enforcing access permissions~\cite{riscvIommu}.\claim{B-SPEC} Translation also introduces a shared service path. A device can reuse its own I/O translation lookaside buffer (IOTLB) entries and still wait behind another device's misses. Prior work has already established IOMMU interference, shared-translation bottlenecks, and mechanisms for hit-under-miss service~\cite{bermuda2025,kurth2018,hero2017}.\claim{L-PREHISTORY} The performance question addressed here is more specific: in a fixed, executable RISC-V implementation, where does the waiting occur, which neighboring operation can extend it, and what does protecting one stream cost the other?

Several closely related studies establish the context. Bermuda measures interference between DMA streams on an Arm SMMUv2 platform, including transfer-size dependence~\cite{bermuda2025}.\claim{L-BERMUDA} Koenig et al. evaluate IOMMU-based shared virtual addressing on a RISC-V field-programmable gate array (FPGA) system, including memory-latency and host-traffic effects~\cite{koenig2025}.\claim{L-DATE} Rodr\'{i}guez's thesis evaluates the same platform family and recognizes contention among its devices~\cite{rodriguez2024}.\claim{L-RODRIGUEZ} Zhao et al. evaluate multiple I/O address-translation caches sharing a translation server, while Hur and Ro evaluate concurrent clients and IOTLB policies using a RISC-V processor configuration and a system memory management unit (SMMU) model~\cite{zhao2025,hur2025}.\claim{L-PROC}\claim{L-HUR} These are direct performance precedents. The contribution of this paper is the depth and traceability of an implementation-specific characterization, rather than discovery of the shared-translation problem.

Aggregate slowdown alone cannot identify a translation bottleneck. A co-running neighbor competes for data movement even if both streams have no measured translation misses. Increasing its address working set changes translation activity, data addresses, and service timing together.\claim{M-CONTROLS} Likewise, a longer internal request lifetime includes waiting and service, and overlapping request lifetimes cannot be added into an elapsed-time breakdown.\claim{M-OBSERVER} We therefore combine a fixed-task co-run control, segmented hardware observations, and a bounded intervention on one identifiable response.

The main measurements use direct wiring of the IOMMU's data-structure request path (DS0). An additional configuration inserts diagnostic delay (DS50) to test sensitivity to that service path.\claim{M-DS} In DS0, growing the aggressor from A1 to A8 and A64 increases median victim hardware elapsed time by \DSZeroAEightElapsedExtra\% and \DSZeroASixtyFourElapsedExtra\%, respectively. The corresponding DS50 increases are \DSFiftyAEightElapsedExtra\% and \DSFiftyASixtyFourElapsedExtra\%.\claim{R-ELAPSED} These comparisons do not assign the difference to pure page-table-walk cost. The DS0 high-pressure conditions also retain small nonzero victim hardware performance monitoring (HPM) miss counts; their software measurement boundary differs from the hardware timing window.\claim{R-MISS}

The paper contributes three connected bodies of evidence. The fixed repaired FPGA baseline quantifies co-run loss under autonomous supply and records measurement windows and source identity.\claim{M-BASELINE}\claim{M-IDENTITY} Historical request-level experiments identify a specific path by which an aggressor's delayed page-table-walk response extends a neighboring hit's wait, separately from capacity replacement.\claim{R-E3}\claim{R-CAPACITY} Historical mitigation experiments account for progress on both sides and preserve a combination regression whose full board-level cause remains unresolved; invalidation scope is a separate simulation supplement.\claim{R-MITIGATION}\claim{R-INVALIDATION} All conclusions are bounded by the tested implementation, load, and evidence level.
